% Chapter 2 -- Background and Literature Review.
% Body only: the \chapter command lives in main_full.tex.
% Source: author's Full report 1.docx (Full Report Draft), 7 Sep 2026.

This chapter introduces the main concepts and prior work on which the
rest of the dissertation is based. It begins with an overview of
reinforcement learning and explains why the portfolio trading problem
can be formulated using the reinforcement learning framework. It then
presents the two algorithm families used in this study: value-based
methods, from Q-learning extending to Deep Q-learning (DQN), and
policy-gradient methods, from REINFORCE to Proximal Policy Optimisation.

The second half of the chapter focuses on reinforcement learning in
financial trading. It reviews findings from previous studies,
particularly how researchers have evaluated trading agents and the
limitations that commonly appear in this area of research.

The chapter also considers the market-efficiency perspective, which
provides an important baseline for measuring trading performance---judging whether a learning algorithm can consistently find profitable
trading opportunities.

The chapter concludes by bringing these points together and identifying
the small gap addressed by this dissertation---whether financial
performance is sufficient to establish that a trained trading policy
actually uses the information available to it.

\section{Reinforcement Learning for Sequential Decisions}
Most machine learning taught at the master\textquotesingle s level is
based on supervised learning, where a model is given inputs together
with the correct answers and learns to reproduce the desired output.
Reinforcement learning (RL) works differently. Instead of being given
the correct action, the learner, known as the \textbf{agent}, interacts
with an environment and learns from the results of its own decisions~\cite{sutton2018}.

This interaction is a repeated process. At each step, the agent observes
the current state, selects an action, and receives a reward, then moves
to the next state. The reward indicates how useful the action was, but
it does not directly tell the agent whether the action was correct. The
agent must therefore learn through experience which actions are likely
to produce higher rewards over time.

Trading fits naturally into this framework. A trader looks at the
markets, decides whether to buy or sell more units of shares, and then
waits for the effect of that decision on his portfolio account. A
decision made early in a month affects the portfolio permanently from
that point onwards. The agent must therefore learn not only whether a
decision was profitable, but also how earlier actions contribute to
later outcomes. This is known as the \textbf{credit-assignment problem}
and is an important part of the mathematical formulation presented in
Chapter 3.

The absence of explicit labelled answers also creates a practical issue.
In supervised learning, an incorrect prediction can be compared directly
with a known target during training. In reinforcement learning, however,
there is no equivalent answer to compare against during training. An
experiment can therefore run without producing an obvious error in the
environment, reward calculation, state representation, or action
handling and still learn the wrong behaviour. This was an important
consideration in this project and is one reason the verification tests
were developed for use before the main experiments, as described in
Chapter 4 and Appendix A.

The broader history of reinforcement learning systems shows this
approach can learn complex sequential tasks. TD-Gammon, for example,
learned to play backgammon strongly through self-play~\cite{tesauro1995}. Later,
Deep Q-learning achieved human-level performance across several Atari
games using raw image inputs~\cite{mnih2015}. The analogy with trading is not
that financial markets are equivalent to games but that in both cases,
the agent is not given the correct action for each situation. Instead,
they learn from a numerical reward or game score. In a trading
environment, financial outcomes such as changes in portfolio wealth play
this role.

The standard mathematical framework for representing sequential problems
is the \textbf{Markov Decision Process (MDP)}. Its foundations are
closely tied to Bellman\textquotesingle s work on dynamic programming~\cite{bellman1957}, and Sutton and Barto present a modern reinforcement learning
formulation~\cite{sutton2018}. An MDP represents a problem through states,
actions, transitions to the next state, rewards, and a discount factor,
assuming that the current state contains all the information needed to
make the next decision. In financial markets, however, it is difficult
to show that any collection of market observations can fully capture the
information relevant to future prices.

For this reason, Chapter 3 uses a practical state representation rather
than claiming that the Markov assumption is perfectly satisfied: it
captures market, portfolio, and timing information rather than claiming
to reconstruct the complete market state.

The study also examines this design experimentally by testing whether
adding drawdown information to the state changes the behaviour or
performance of the learned agents. The next section discusses the two
main families of reinforcement learning algorithms considered in this
dissertation.

\section{Two Families of Learning Algorithms}
The algorithms used in this dissertation belong to two different
families of reinforcement learning methods. This section discusses the
main ideas behind each family and explains why comparing them on the
same trading problem is useful.

Value-based methods estimate the value of taking an action and derive a
policy from those estimates, whereas policy-gradient methods select the
action with the highest probability to optimise the policy itself. These
differences matter because they offer two different learning mechanisms
for the same trading problem.

\subsection{Value-based methods: from Q-learning to DQN}
The Value-based family approaches the question of what action to take
indirectly. Rather than learning the policy directly, the agent learns
how valuable different actions are in a given state and then selects the
action with the highest estimated value.

This approach started from \textbf{temporal-difference learning}~\cite{sutton1988}, which updates value estimates using the difference between
sequential predictions rather than waiting until the end of an episode
to observe the final outcome. \textbf{Q-learning} turned this idea into
an off-policy method that directly estimates the action-value function~\cite{watkins1992}. Under appropriate conditions, including continued exploration
of state-action pairs and suitable learning-rate decay, tabular
Q-learning converges to the optimal action values for finite state
spaces. This convergence property does not depend on the particular
actions used while learning.

In tabular Q-learning, a separate value is stored for every state-action
pair. This works well when the number of possible states is small, such
as in a simple grid-world, but becomes impractical when the state
contains continuous financial information. Two very similar market
states would also be treated as completely separate entries in the
table. As a result, experience gained in one state would not
automatically help the agent in a similar state.

\textbf{Deep Q-learning (DQN)}~\cite{mnih2015} addresses this problem by
replacing the table with a neural network that estimates the action
values. The network can therefore generalise across similar states,
allowing information to be shared, but it introduces instability.
However, combining neural networks with bootstrapping and off-policy
learning can make training unstable~\cite{tsitsiklis1997, sutton2018}.

DQN then introduces two important mechanisms to improve this stability.
\textbf{Experience replay}~\cite{lin1992} stores previous transitions in a
replay buffer and samples them later in random minibatches during
training. This reduces the strong correlation present in consecutive
experiences, allowing each transition to contribute to multiple updates.
A \textbf{target network} is also used as a periodically updated copy of
the value network to construct the learning target, preventing it from
changing with every network update as the parameters are being
optimised.

These mechanisms are highly useful for financial data, where consecutive
daily observations are naturally correlated. Both were implemented from
scratch in this project, as described in Chapter 4. As a result, the
implementation shows the underlying learning mechanics rather than
treating the algorithm as a library component.

Several Other improvements have since been proposed to address the
limitations of the original DQN formulation. For example, \textbf{Double
Q-learning}~\cite{vanhasselt2016} was developed partly to reduce overestimation bias
in action values that can arise in standard Q-learning when the same
estimates are used to select and evaluate actions.

Such modifications are relevant to the broader DQN literature or
discussed as possible future improvements but are not included in this
study's present implementation, allowing the core DQN method to remain
relatively simple and easier to verify.

\subsection{Policy-gradient methods: from REINFORCE to PPO}
Policy-gradient methods take a more direct approach to the decision
problem. Instead of first estimating the value of each action to derive
a policy from those estimates, the agent learns a policy that maps
states to actions and adjusts its parameters to increase the expected
return. The policy is represented by a differentiable model, typically a
neural network, whose parameters are updated for the objective.

This study uses the foundational policy-gradient algorithm
\textbf{REINFORCE}~\cite{williams1992}. One of its main advantages is its
simplicity, as it makes the relationship between the policy, observed
states, selected actions, and expected return relatively transparent. It
does not require bootstrapping, experience replay, or a separate target
network. The policy is updated using returns from complete episodes,
making the learning process straightforward to understand. For these
reasons, REINFORCE is used as the main policy-gradient implementation in
this dissertation.

The main weakness of REINFORCE is the high variance of its updates. The
learning signal is based on Monte Carlo returns, which can vary
considerably between episodes. This can make learning slow and unstable
when rewards are noisy.

Later policy-gradient methods introduced techniques to address this
problem. Actor-critic methods use a learned value estimate as a baseline
to reduce the variance of the policy update~\cite{konda2000}. The actor
represents the policy while the critic estimates value, allowing the
learning signal to use an advantage rather than the raw return\textbf{.
A2C} is one example of this approach~\cite{mnih2016}. \textbf{Generalised
Advantage Estimation (GAE)}~\cite{schulman2016gae} is another way of controlling the
bias-variance in the advantage estimate.

A second problem is the size of policy updates. An update that changes
the policy too aggressively can cause training to become unstable and
degrade performance sharply. Trust-region methods~\cite{schulman2015trpo} address this
by restricting how far the policy can change between successive updates.
\textbf{Proximal Policy Optimisation (PPO)}~\cite{schulman2017ppo} uses a simpler
approach by clipping the policy update when it moves too far from the
previous policy.

PPO is widely used in reinforcement learning applications, including
trading frameworks such as FinRL~\cite{liu2020finrl}, and is therefore included in
this study as a library-based policy gradient method for comparison
while implementing REINFORCE from scratch. This allows for both a
transparent baseline policy-gradient method and a more modern,
stabilised method

The trading environment also contains actions that are not always valid.
For example, an agent should not be allowed to sell when it has no
position or buy when there is insufficient cash. All three learning
methods use \textbf{invalid-action masking} to prevent such actions from
being selected while preserving only valid actions. The use of action
masking in reinforcement learning, including its interaction with
policy-gradient methods, has been examined by Huang and Ontañón~\cite{huang2022masking}. In this dissertation, masking is part of the common
environment rather than a difference between algorithms, ensuring that
the comparison is made under the same set of trading constraints.

\subsection{Why compare the two families?}
The two algorithm families have different implementation trade-offs.

Value-based methods such as DQN can learn from individual transitions
and reuse previously collected experience through replay. However, they
rely on bootstrapping value estimates, meaning the agent learns partly
from estimates produced by other learned estimates in the same learning
process.

Policy-gradient methods instead optimise the policy directly with
respect to the objective. In its basic form, REINFORCE provides a
balanced policy-gradient estimate, but it requires new experiences for
each update from fresh trajectories, and gradient estimates can suffer
from high variance. PPO improves the stability of this family through
constrained policy updates.

Neither family can therefore be assumed to be better for the portfolio
trading problem considered here. Comparing them in the same environment
and under the same evaluation procedure is therefore useful. If both
families develop similar behaviour, this may suggest that the behaviour
is mainly caused by the structure of the problem, reward, state
representation, or environment. If their behaviour differs, the
difference becomes evidence that the choice of learning method
influences the resulting strategy.

This comparison becomes an important part of the experimental design in
Chapter 5, as it is not limited to only asking which algorithm produces
the highest financial return. It also asks whether the resulting
policies learn fundamentally different ways of responding to the trading
environment. The next section discusses how previous research applies
these algorithms to financial markets.

\section{Reinforcement Learning for Trading}
The use of reinforcement learning (RL) for trading has been studied for
many years, but the application comes with several well-known
challenges. One of such early influential work is Moody and Saffell on
direct reinforcement methods~\cite{moody2001}, which uses gradient-based methods
to learn trading policies directly. Their work raised an important
point: the reward should account for risk. If an agent is rewarded only
for increasing profit, it may have a strong incentive to remain highly
exposed to the market without explicitly rewarding control of downside
risk.

This observation is highly relevant to this dissertation, giving us
reason to run two experiments in Chapter 5: first, to examine what
happens when the agents are trained using a simple wealth-change reward
and the objective is tied only to portfolio growth; and second, whether
adding a drawdown penalty to the reward changes their behaviour.

Drawdown is also introduced separately as an additional state feature.
These two uses of drawdown act differently. Adding drawdown to the state
changes the information available to the agent, whereas adding a
drawdown penalty in the reward changes what the agent is encouraged to
optimise. One changes what the agent can see, while the other changes
what it is rewarded for. However, the specific drawdown penalty used in
this dissertation is also an experimental design choice to be evaluated
rather than a reward formulation from previous literature.

Keeping these changes separate allows their effects to be examined
separately. The experiments determine whether either change affects the
learned behaviour, rather than assuming what the resulting behaviour
should be.

Later work extended these ideas using more complex models. Deng et al.~\cite{deng2017}, for example, combined direct reinforcement with deep recurrent
representations. Jiang et al.~\cite{jiang2017} applied deep reinforcement
learning to multi-asset portfolio management using cryptocurrency data,
using a different formulation for the action, which represents portfolio
allocation rather than a discrete buy, sell, or hold decision.

More recently, frameworks such as FinRL~\cite{liu2020finrl} have made it easier to
build trading environments and apply these reinforcement learning
algorithms to financial experiments, but it often defaults to PPO-based
methods.

These frameworks are useful for experimentation, but they can also hide
parts of the underlying learning process. This was one reason for
implementing REINFORCE and DQN from scratch in this dissertation,
allowing their behaviour and training mechanisms to be examined
directly.

The study most closely related to the present work is Théate and Ernst~\cite{theate2021}. They evaluate a DQN-based trading agent on a test portfolio
set of 30 stocks, comparing it with buy-and-hold and traditional active
strategies. Their results report that the DQN performs better on average
than the benchmark strategies, although they include several active
strategies rather than only passive buy-and-hold, which is more relevant
to this dissertation. For SPY, which is also used in this dissertation,
their reported Sharpe ratio for the DQN was 0.834, the same as
buy-and-hold.

Across multiple stocks, they describe the agent\textquotesingle s
performance as identical or very close to passive strategies and suggest
that the agent can tend towards a passive strategy when active trading
becomes more uncertain. The variations between training runs showed that
the agent traded less as transaction costs increased, eventually
becoming passive. They also report non-negligible variance across
identical training runs. These findings make Théate and Ernst study
relevant to this study's evaluation.

There are, however, important differences between their experimental
design and the one used here. Their study trains a separate agent for
each stock, allows short positions, and uses a continuous test period,
whereas this dissertation studies one ETF, does not permit shorting, and
uses independent monthly episodes with a fixed trade size. Their
benchmarks also differ from this study's buy-and-hold strategies,
preventing a direct numerical comparison between the studies.

More importantly, although their results show that the learned agent can
become increasingly passive, they do not directly test how the learned
policy generates its decisions: whether the policy responds to its
observations or simply follows a fixed trading pattern. Both behaviours
could produce similar trading trajectories.

This difference is fundamental to this dissertation, as a policy can
achieve a strong return without making meaningful use of market
information, particularly when the asset always goes up over time.
Chapter 5 therefore includes a specific behavioural analysis to
determine whether the actions of each trained agent actually change when
its observed state changes while the available actions remain the same.

Read together, the previous studies reviewed here vary widely in how
RL-based trading is formulated and evaluated. Some focus on individual
assets (Moody and Saffell~\cite{moody2001}, Deng et al.~\cite{deng2017}, and Théate and
Ernst~\cite{theate2021}), while others use multi-asset portfolios (Jiang et al.~\cite{jiang2017} and Guan and Liu~\cite{guan2021xrl}). Action representations range from
discrete trading decisions~\cite{moody2001, deng2017, theate2021} to continuous portfolio
weights~\cite{jiang2017}. Reward functions also vary, from direct profit or
return objectives~\cite{deng2017, jiang2017, theate2021} to risk-adjusted objectives such as
Moody and Saffell\textquotesingle s differential Sharpe formulation~\cite{moody2001}. Most papers include transaction costs, but their assumptions
and representations differ. Benchmark strategies also vary, from
traditional trading strategies~\cite{theate2021} to portfolio-selection methods~\cite{jiang2017}.

There is also variation in the results across seed runs. Some studies
acknowledge differences between training runs~\cite{theate2021}, but reported
results are often based on individual runs. None of the reviewed studies
also describe a procedure in which all experimental design decisions are
selected using a validation period before being used on an unseen test
period. Behavioural analysis is similarly limited, with previous work
using methods such as trading trajectories~\cite{theate2021} or feature
attribution~\cite{guan2021xrl}.

These differences make direct comparison between studies difficult and
highlight an important issue: the evaluation procedure can be just as
important as the choice of algorithm. The studies reviewed in this
section were selected because they are particularly relevant to this
dissertation's design, especially their use of RL algorithms,
transaction costs, and benchmarks. Fischer\textquotesingle s survey~\cite{fischer2018} provides a broader review of RL in financial markets.

The next section therefore focuses more specifically on how to evaluate
trading agents and on the weaknesses that can arise when financial
performance is considered without sufficient behavioural and
experimental controls.

\section{Evaluating Trading Agents Honestly}
A review of reinforcement learning in finance~\cite{fischer2018}, together with
research on financial machine learning~\cite{bailey2014, lopezdeprado2018}, identifies several
common problems that can make trading results appear stronger than they
are. Four of these issues are relevant to this dissertation, and each
one influenced the experiment design. These are not claims that every
reviewed study states as problems; they are risks that an experiment
should account for when being designed.

The first issue is \textbf{weak benchmarks}. A learned trading strategy
can appear successful if it is compared with a benchmark that invests
less capital or does not represent a realistic alternative. For this reason,
this study uses buy-and-hold as the main benchmark, with the same
initial capital as the learned agents. The trade size was also
calibrated in Chapter 5 so the agents could reach a comparable level of
market exposure. This keeps the comparison fair, as the agents are not
disadvantaged simply by being unable to invest enough capital.

The second issue is \textbf{look-ahead bias.} Financial data are ordered
in time, so treating it like randomly ordered observations and shuffling
the data before splitting could place future information in the training
set and produce overly optimistic results. The data in this study are
therefore divided chronologically: 2018--2022 for training, 2023 for
validation, and 2024--2025 as the held-out test period. The test months
were not used to select hyperparameters or make any design decisions, as
described in Chapter 4.

The third issue is \textbf{ignoring transaction costs.} A strategy that
trades frequently may look profitable when every trade is assumed to be
free, but the same strategy could lose much of its return once real
trading costs are included. In this dissertation, transaction fees are
part of the trading environment and are therefore included while the
agents are learning rather than only after training. Chapter 5 also
varies the assumed fee to examine whether the learned strategies trade
less when trading becomes more expensive.

The fourth issue is \textbf{backtest overfitting.} Repeatedly adjusting
a strategy after looking at its test results gradually turns the test
set into another training set. Bailey et al.~\cite{bailey2014} discussed this
problem for investment backtests, and López de Prado~\cite{lopezdeprado2018} discussed
methodology for controlling leakage and data reuse. The defence here is
that all design choices were selected using the training and validation
data, then evaluated once on the held-out test months, without further tuning.

A related concern comes from deep reinforcement learning itself.
Henderson et al.~\cite{henderson2018} show that reported \textbf{results can vary
substantially across random seeds}, making a single successful training
run a weak basis for comparison. Each configuration in this study is
therefore trained with six different seeds, and the spread is reported
alongside the mean rather than presenting one favourable run as
representative.

Financial economics provides more context for understanding the results.
Under the \textbf{efficient market hypothesis,} current prices should
already reflect the information contained in past prices~\cite{fama1970}. If
this assumption holds, an agent whose market information consists only
of price history should not be expected to consistently earn higher
risk-adjusted returns than the market itself. In this study,
buy-and-hold on the single fund provides the practical benchmark against
which this expectation is tested.

Portfolio theory also emphasises that return cannot be considered
separately from risk~\cite{markowitz1952}. A strategy that earns less may still be
useful if it achieves a large enough reduction in risk. The evaluation
therefore reports mean return together with drawdown and the Sharpe
ratio~\cite{sharpe1994}. No single metric decides whether a strategy is
successful; the results are read together so that the trade-off between
return and risk remains visible.

This approach does not aim to guarantee a positive result for the
learning algorithms. If an agent does not beat buy-and-hold, that
outcome is treated as a finding to explain rather than a failure to
hide. Understanding why a strategy does not improve on a simple
benchmark is therefore just as important as reporting successful
results. However, even a carefully controlled financial experiment only
establishes what happened to the portfolio; it does not establish what
the agent learned to do. That distinction leads to the gap considered in
the next section.

\section{The Gap Addressed In This Study}
Financial performance alone does not tell us what a trading agent has
actually learned. An agent can follow a simple rule, such as buying
whenever possible, and still perform well during a rising market without
using any market information. At the same time, an agent that genuinely
responds to market conditions may earn less over the same period. For
this reason, financial performance and agent behaviour need to
be evaluated separately, and the second is needed to interpret the
first.

This distinction matters because most of the trading studies reviewed in
this chapter focus mainly on financial measures such as return, profit,
Sharpe ratio, and drawdown. More recent work also examines how trained
models make their decisions, including feature-attribution methods by
Guan and Liu~\cite{guan2021xrl}. These methods can identify which inputs are
related to a model\textquotesingle s output. However, they do not
establish whether the policy changes its actions as its inputs change. A
policy that produces the same action regardless of its observations
leaves no variation in action for an attribution method to explain.

Checking whether a learned agent responds to its observations is already
used in reinforcement-learning research and in the wider evaluation of
machine-learning systems. This dissertation therefore does not claim to
introduce sensitivity testing or behavioural analysis as new concepts,
as there are related ideas that exist in input sensitivity.

The more specific issue addressed here is how to test dependence on
state information in a trading environment with constraints on actions.
The agent\textquotesingle s cash and holdings determine which actions
are legal: it cannot buy when it has insufficient cash and cannot sell
when it has no position. The set of available actions can therefore
change during an episode, and the selected action may change simply
because an action became unavailable rather than because the agent
responded to market information, and this would risk counting these
forced changes as evidence of state use.

To account for this, this dissertation defines state dependence
intentionally: \textbf{trained policies} \textbf{must select different
actions in situations where the same actions are legally available, with
the difference occurring in the observed state}. Conditioning the
analysis on action masking reduces the chance of misinterpreting changes
caused by the constraints as evidence that the agent is using its
observations.

The diagnostic, however, answers a limited question --- does the policy
show any evidence of changing its decisions in response to its observed
state? It does not identify which feature caused the change, whether the
response was economically sensible, or whether it improved financial
performance.

Among the trading studies reviewed here, policy responsiveness is not
reported as a quantitative evaluation result alongside financial
performance. This is a narrower claim than saying previous research has
never used sensitivity testing or behavioural analysis. Established
practices such as chronological splits, transaction-cost modelling,
benchmark comparisons, and multi-seed reporting are also not presented
as new contributions here.

Théate and Ernst~\cite{theate2021} show why this distinction matters. Their DQN
moved towards passive trading under some conditions, but passive
behaviour could come from a policy that responds to market information
or from one that follows a fixed trading rule. Without directly testing
the relationship between the observed state and the chosen action, both
cannot be explained separately.

This dissertation addresses this gap through a controlled comparison of
three different learning algorithms. REINFORCE, Deep Q-learning, and PPO
are trained using the same observed state, the same masked action space,
the same fee-inclusive reward, and the same chronological splits. Their
financial performance is compared against benchmark strategies with the
same capital, and their behaviour is examined using state-dependence
analysis. The diagnostic is reported for each seed alongside the
financial results rather than being treated as a technical
implementation check.

Two further experiments examine whether changing the information
available to the agent or changing its incentive affects the learned
behaviour. The first adds drawdown to the state, allowing the agent to
observe its recent portfolio losses. The second adds a drawdown penalty
to the reward, changing the incentive to take risks. One changes what
the agent can see, while the other changes what the agent is encouraged
to optimise.

The contribution of this dissertation is therefore not the invention of
a new evaluation method, but the controlled combination of experimental
practices established from prior works to measure the
agent\textquotesingle s state-dependent behaviour along with financial
performance. This allows the evaluation to distinguish between an agent
that earns a strong return by following a simple fixed rule and one
whose decisions actually respond to the information provided.

These points lead directly to the four research questions introduced in
Chapter 1. The first asks whether the learned agents can outperform
meaningful benchmarks after transaction costs. The second asks whether
their decisions depend on the information in their states. The third
asks whether adding drawdown to the state changes behaviour or
performance. The fourth considers whether penalising drawdown through
the reward produces a different risk-sensitive outcome. The next chapter
formalises these questions through the MDP formulation, state
representation, trading rules, reward functions, and the algorithms used
in the experiments.
